\documentclass[11pt]{article}
\usepackage[margin=1in]{geometry}
\usepackage{amsmath,amssymb,amsthm,mathtools}
\usepackage[expansion=false]{microtype}
\usepackage[hidelinks]{hyperref}
\usepackage{xcolor}
\providecommand{\doi}[1]{\href{https://doi.org/#1}{doi:#1}}

\newtheorem{theorem}{Theorem}
\newtheorem{lemma}[theorem]{Lemma}
\newtheorem{proposition}[theorem]{Proposition}
\newtheorem{corollary}[theorem]{Corollary}
\theoremstyle{definition}
\newtheorem{definition}[theorem]{Definition}
\newtheorem{openproblem}[theorem]{Open Problem}
\theoremstyle{remark}
\newtheorem{remark}[theorem]{Remark}

\newcommand{\Z}{\mathbb{Z}}
\newcommand{\C}{\mathbb{C}}
\newcommand{\ord}{\operatorname{ord}}
\newcommand{\Spec}{\operatorname{Spec}}

\hypersetup{
  pdftitle={Exact Spectral Inheritance in Prime-Power Resolution Towers},
  pdfauthor={Elias De Jes\'us},
  pdfkeywords={transfer operator, spectral gap, Collatz, prime-power tower, Wieferich prime, equidistribution}
}

\title{\textbf{Exact Spectral Inheritance in\\ Prime-Power Resolution Towers}\\[2pt]
\large A Splitting Theorem for Affine-Cocycle Transfer Operators,\\
with Application to the Collatz Carry Sum}
\author{Elias De Jes\'us\\[2pt]
\small Independent Researcher\\
\small ORCID: \href{https://orcid.org/0009-0007-0190-9143}{0009-0007-0190-9143}\quad
\href{mailto:dejesuselias10@gmail.com}{dejesuselias10@gmail.com}}
\date{Technical note, Version 1.0 --- October 2026}

\begin{document}
\maketitle

\begin{abstract}
For a prime power $q=p^e$, the accelerated Collatz carry sum generates a multiplicative walk on $K_e=\langle 2,3\rangle\le(\Z/q\Z)^\times$ and an associated family of phase-twisted transfer operators $L_{e,t,x}$, whose dominant twisted eigenvalue controls the rate of equidistribution of carry-sum residues at fixed $p$~\cite{DJ-fixedprime}. This note settles how this spectral picture behaves as the resolution $e$ is refined, by proving an exact structural theorem for the whole tower $(L_e)_{e\ge1}$. We show that the natural pullback $J_e:H_e\to H_{e+1}$ along the reduction $\pi_e:K_{e+1}\to K_e$ intertwines the operators ($L_{e+1}J_e=J_eL_e$) at \emph{every} prime-power step, with no hypothesis on multiplicative orders, and that its image is a \emph{reducing}, not merely invariant, subspace: the orthogonal complement $Z_{e+1}$ (functions of zero sum on every fiber of $\pi_e$) is also exactly invariant, giving an orthogonal direct-sum splitting $L_{e+1}\cong L_e\oplus N_{e+1}$. This is strictly stronger than the block-upper-triangular structure a bare invariant subspace would give, and it has an immediate consequence: the spectral-gap quantity $g_0(p,e)$ governing equidistribution is non-increasing in $e$ from $e=1$, for every prime $p\nmid 6$, and hence convergent. This resolves, in the monotone-and-convergent direction, an open question raised by preliminary numerical data that had suggested an ongoing decay. The key to universality is to separate two quantities that coincide generically: the growth $c_e=r_{e+1}/r_e$ of the order of the resolution base $b=2$, and the fiber size $m_e=|K_{e+1}|/|K_e|$. At Wieferich-type steps the order of $2$ stalls ($c_e=1$) while the fiber may still grow ($m_e=p$); the coarse matrix entry is then realized by a single fine increment per target instead of $p$ increments combining geometrically, and inheritance survives. The base-2 Wieferich prime $1093$ is the concrete instance in the Collatz case, and the base-3 Wieferich prime $11$ is an anomaly of a different channel; neither affects inheritance. The theorem and its proof are stated for a general affine-cocycle operator family; the Collatz carry sum is the motivating instance, not a hypothesis.
\end{abstract}

\medskip
\noindent\textbf{Keywords:} transfer operators; spectral gap; prime-power towers; reducing subspaces; Wieferich primes; accelerated Collatz map; equidistribution.\\
\textbf{MSC 2020:} 11B83 (primary); 37A30, 47A15, 11A07 (secondary).

\bigskip
\noindent\fbox{\parbox{\dimexpr\linewidth-2\fboxsep-2\fboxrule}{\small
\textbf{Scope of claims.}
\emph{Proved here, generically} (Theorems~\ref{thm:intertwine}, \ref{thm:split}, \ref{thm:mono}): the intertwining identity; the exact orthogonal splitting; monotonicity and convergence of $g_0(p,e)$ at every level $e\ge1$, for every prime $p\nmid ab$, with no order-lifting hypothesis.
\emph{Proved here, number-theoretically} (Propositions~\ref{prop:wieferich}, \ref{prop:lte}): the equivalence of the Fermat-quotient and order-stall forms of the Wieferich condition, and the exact order-lifting schedule, which governs the shape of the increment ladder (not the validity of the splitting).
\emph{Verified computationally} (exact/high-precision, small cases): the intertwining identities in exact rational arithmetic over 3{,}134 prime-power steps (including 262 steps where the order of $b$ stalls) and at $p=1093$; the splitting and eigenvalue-union formula numerically, for several primes and non-Collatz controls.
\emph{Explicitly open:} whether a non-primitive stratum's new-sector operator can ever exceed its own inherited eigenvalue (Open Problem~\ref{op:local}); whether $\lim_e g_0(p,e)$ is strictly positive.
\emph{Not claimed:} anything about Collatz cycle closure, divergence, or the Collatz conjecture. This note is purely about the statistical/spectral structure of the carry-sum residue system.}}

\section{Introduction}

The $3x+1$ problem and its accelerated (Syracuse) form have a long literature on the statistical behaviour of iterates and their residues; see Lagarias~\cite{Lagarias1985,Lagarias2010} for surveys and Tao~\cite{Tao2022} for the strongest current results on equidistribution of Syracuse random variables modulo powers of $3$.

\emph{Fixed-Prime Equidistribution of Accelerated Collatz Carry Sums}~\cite{DJ-fixedprime} proves that, for the accelerated Collatz carry sum $C_L(D)=\sum_{j<L}3^{L-1-j}2^{s_j}$ over valuation words $D$, the residue $C_L(D)\bmod p$ equidistributes as $L\to\infty$, for every fixed prime $p\nmid 6$, at an explicit geometric rate governed by a transfer-operator spectral gap $g_0(p)$ (transfer operators in this sense are the standard tool of thermodynamic formalism; see Baladi~\cite{Baladi2000}). The companion instrument \texttt{amgap}~\cite{DJ-amgap} extends the construction to prime powers $q=p^e$ and reports two data points suggesting the gap might shrink with $e$: $g_0(11,1)=0.3340$, $g_0(11,2)=0.3225$, explicitly flagging whether $g_0(p,e)\to0$ as $e\to\infty$ as open.

A preliminary numerical extension (unpublished working notes) computed $g_0(11,3)\approx0.305$, apparently continuing the downward trend, and an initial structural analysis of the prime-power tower proposed a recursion $B_p(e)=\max\{P_p(e),B_p(e-1)\}$ (with $P_p(e)$ the dominant primitive-stratum eigenvalue and $B_p(e)$ the overall gap-defining quantity), based on an embedding of lower-level spectra into higher levels. On audit, that recursion was found to rest on an unproved assumption, namely that an embedded (inherited) eigenvalue remains globally dominant within its own stratum, and the monotonicity conclusion was downgraded to conditional.

This note resolves the question. The missing structural fact is that the pullback along the reduction map is not merely an invariant-subspace embedding but splits the operator as an honest orthogonal direct sum. This upgrades the available inequality to an exact formula, from which monotonicity follows by a short, elementary argument (Theorem~\ref{thm:mono}). No arithmetic hypothesis is needed: an earlier version of the argument assumed that the order of the resolution base $b$ grows by exactly $p$ at each step, but the geometric weighting absorbs order stalls exactly (Lemma~\ref{lem:ladder}), so the splitting holds at every step, including at Wieferich-type primes (Section~\ref{sec:arith}).

In the language of the Closure--Separation--Stability audit protocol for reduced representations~\cite{DJ-css}, the reduction $\pi_e$ is a coarse-graining of the residue state space. Theorem~\ref{thm:split} says that this coarse-graining is exact in a strong sense: the coarse observables close under the dynamics (the level-$e$ operator is recovered exactly), and the information discarded by the reduction separates into its own invariant sector rather than leaking back into the coarse one.

\section{The operator family}

\begin{definition}[Affine-cocycle transfer operator]\label{def:op}
Fix an odd prime $p$ and coprime integers $a,b\ge2$ with $p\nmid ab$. For $e\ge1$, let $q=p^e$ and $K_e=\langle a,b\rangle\le(\Z/q\Z)^\times$, a cyclic group (every subgroup of $(\Z/p^e\Z)^\times$ is cyclic for odd $p$~\cite{IrelandRosen}). Let $r_e=\ord_q(b)$. For $t\in\Z/q\Z$ and $x$ in the open unit disk, define $L_{e,t,x}:H_e\to H_e$ on $H_e:=\C^{K_e}$ by
\[
(L_{e,t,x}f)(u')=\sum_{u\in K_e} e_q(tu)\,\frac{x^{d_0(u,u')}}{1-x^{r_e}}\,f(u),
\]
where $e_q(y)=\exp(2\pi i y/q)$ and $d_0(u,u')\in[1,r_e]$ is the least $d$ with $u'\equiv u\,b^d a^{-1}\pmod q$ (the entry is $0$ if no such $d$ exists).
\end{definition}

The Collatz instance is $(a,b)=(3,2)$, with $K_e=\langle 2,3\rangle$ the reachable states of the multiplicative walk $u_{j+1}=u_j b^{D_j}a^{-1}$ underlying $C_L(D)\bmod q$~\cite{DJ-fixedprime}. The relevant evaluation point is the saddle $x_*=1-1/\log_2 a$ (for $(a,b)=(3,2)$: $x_*=1-1/\log_2 3$). At $t=0$, for each source $u$ the targets $u\,b^da^{-1}$, $d\in[1,r_e]$, are distinct, so every column of the (nonnegative, for real $x\in(0,1)$) matrix sums to $\sum_{d=1}^{r_e}x^d/(1-x^{r_e})=x/(1-x)$; hence $\rho(L_{e,0,x})=x/(1-x)$ and $\rho(L_{e,0,x_*})=\log_2 a-1$ exactly, independent of $p$ and $e$~\cite{DJ-fixedprime}.

We fix the normalized inner product $\langle f,g\rangle_e:=\frac{1}{|K_e|}\sum_{u\in K_e}f(u)\overline{g(u)}$ throughout (uniform probability measure on $K_e$), the normalization making the resolution maps below isometric.

\begin{definition}[Resolution maps]
$\pi_e:K_{e+1}\twoheadrightarrow K_e$ is reduction mod $p^e$, restricted to $K_{e+1}$; it is a surjective group homomorphism (generators reduce to generators), so all fibers have the same size
\[
m_e:=|K_{e+1}|/|K_e|=|\ker\pi_e|\in\{1,p\}.
\]
Separately, write $B_e:=\langle b\rangle\le K_e$, so $|B_e|=r_e$, and
\[
c_e:=r_{e+1}/r_e .
\]
Define
\begin{align*}
J_e&:H_e\to H_{e+1}, & J_ef&:=f\circ\pi_e && (\text{pullback}),\\
S_e&:H_{e+1}\to H_e, & (S_ef)(\bar u)&:=\textstyle\sum_{u\in\pi_e^{-1}(\bar u)}f(u) && (\text{fiber sum}).
\end{align*}
\end{definition}

Since each fiber has $m_e$ elements, $\langle J_ef,J_eg\rangle_{e+1}=\frac{m_e}{|K_{e+1}|}\sum_{\bar u}f(\bar u)\overline{g(\bar u)}=\langle f,g\rangle_e$, so $J_e$ is an isometry (in particular injective), and $S_e\circ J_e=m_e\cdot\mathrm{id}$.

\begin{lemma}[Ladder versus fiber]\label{lem:cm}
$c_e=|B_{e+1}\cap\ker\pi_e|$; in particular $c_e\mid m_e$ and $c_e\in\{1,p\}$. All three combinations $(c_e,m_e)=(p,p),(1,p),(1,1)$ occur, and $(p,1)$ does not.
\end{lemma}

\begin{proof}
$\pi_e$ maps $B_{e+1}$ onto $B_e$ with kernel $B_{e+1}\cap\ker\pi_e$, so $r_{e+1}=r_e\,|B_{e+1}\cap\ker\pi_e|$, a subgroup of $\ker\pi_e$. Since $\ker\pi_e$ lies in the order-$p$ kernel of $(\Z/p^{e+1}\Z)^\times\to(\Z/p^e\Z)^\times$, both numbers lie in $\{1,p\}$ and $(p,1)$ is impossible. Examples: $(3,2)$ at $p=13$, $e=1$ gives $(p,p)$; $(3,2)$ at $p=1093$, $e=1$ gives $(1,p)$ (Section~\ref{sec:arith}); $(a,b)=(7,18)$ at $p=5$, $e=1$ gives $(1,1)$, since $7^4\equiv18^4\equiv1\pmod{25}$.
\end{proof}

\begin{lemma}[Increment ladder]\label{lem:ladder}
Fix $u'\in K_{e+1}$ and $\bar u\in K_e$, and let $A(u',\bar u)\subseteq\pi_e^{-1}(\bar u)$ be the set of fine sources joined to $u'$ by some increment. Then $A(u',\bar u)=u'aB_{e+1}\cap\pi_e^{-1}(\bar u)$ is empty precisely when the coarse entry from $\bar u$ to $\pi_eu'$ vanishes; otherwise $|A(u',\bar u)|=c_e$, the fine increments of its elements are exactly $d_0+kr_e$ for $k=0,\dots,c_e-1$, where $d_0=d_0(\bar u,\pi_eu')$, and
\[
\sum_{u\in A(u',\bar u)}\frac{x^{d_0(u,u')}}{1-x^{r_{e+1}}}
=\sum_{k=0}^{c_e-1}\frac{x^{d_0+kr_e}}{1-x^{c_er_e}}
=\frac{x^{d_0}}{1-x^{r_e}} .
\]
The same holds with the roles of source and target exchanged: for fixed $u\in K_{e+1}$ and $\bar u'\in K_e$, the fine targets in $\pi_e^{-1}(\bar u')$ reached from $u$ form $ua^{-1}B_{e+1}\cap\pi_e^{-1}(\bar u')$, of size $0$ or $c_e$, with the same increments and the same sum.
\end{lemma}

\begin{proof}
$u$ is joined to $u'$ iff $u=u'ab^{-d}$ for some $d\in[1,r_{e+1}]$, and distinct $d$ give distinct $u$; so the reachable sources form the coset $u'aB_{e+1}$, which meets a fiber of $\pi_e$ in a coset of $B_{e+1}\cap\ker\pi_e$ or not at all, i.e.\ in $c_e$ points or none (Lemma~\ref{lem:cm}). The source $u'ab^{-d}$ lies over $\bar u$ iff $b^{d}\equiv b^{d_0}\pmod{p^e}$, i.e.\ $d\equiv d_0\pmod{r_e}$, and $[1,c_er_e]$ contains exactly the $c_e$ values $d_0+kr_e$. The sum is a finite geometric series: $x^{d_0}(1-x^{c_er_e})/\big((1-x^{r_e})(1-x^{c_er_e})\big)$. The target version is identical with $u'=ub^da^{-1}$.
\end{proof}

\begin{remark}[Stalled steps]\label{rem:stall}
When $c_e=p$ the coarse weight of an entry is reproduced by $p$ fine increments landing on the \emph{same} target and combining geometrically. When $c_e=1<m_e$, each target receives a single fine increment from a given fiber, with the identical weight $x^{d_0}/(1-x^{r_e})$ (the denominator does not change, since $r_{e+1}=r_e$); the remaining $m_e-1$ fiber members lie in other cosets of $\langle b\rangle$ and have zero entry toward that target, feeding instead the other targets $u'z$, $z\in\ker\pi_e\setminus\{1\}$, of the same fiber. Either way, every coarse entry is reproduced exactly.
\end{remark}

\section{Intertwining and the splitting theorem}

\begin{theorem}[Pullback intertwining]\label{thm:intertwine}
For every $e\ge1$ and every $t\in\Z/p^e\Z$,
\[
L_{e+1,pt,x}\,J_e=J_e\,L_{e,t,x}.
\]
\end{theorem}

\begin{proof}
Fix a target $u'\in K_{e+1}$. Grouping the fine sources by the fibers of $\pi_e$, and using that the phase $e_{p^{e+1}}(ptu)=e_{p^e}(t\bar u)$ is constant on the fiber over $\bar u$, Lemma~\ref{lem:ladder} gives
\begin{align*}
(L_{e+1,pt,x}J_ef)(u')&=\sum_{\bar u\in K_e}e_{p^e}(t\bar u)\,f(\bar u)\!\!\sum_{u\in A(u',\bar u)}\!\!\frac{x^{d_0(u,u')}}{1-x^{r_{e+1}}}\\
&=\sum_{\bar u\in K_e}e_{p^e}(t\bar u)\,\frac{x^{d_0(\bar u,\pi_eu')}}{1-x^{r_e}}\,f(\bar u)
=(L_{e,t,x}f)(\pi_eu'),
\end{align*}
which is $(J_eL_{e,t,x}f)(u')$. No relation between $c_e$ and $m_e$ is used.
\end{proof}

\begin{remark}
The proof never uses $\gcd(t,p)=1$; it lifts any residue $t$ at the lower level by one factor of $p$, not only primitive ones.
\end{remark}

\begin{lemma}[Pushforward intertwining]\label{lem:push}
For every $e\ge1$ and $t\in\Z/p^e\Z$, $S_e\circ L_{e+1,pt,x}=L_{e,t,x}\circ S_e$.
\end{lemma}

\begin{proof}
Proved independently of Theorem~\ref{thm:intertwine}: by the target form of Lemma~\ref{lem:ladder}, with the source $u\in K_{e+1}$ fixed and the sum taken over targets $u'\in\pi_e^{-1}(\bar u')$, the phase being that of the source, $\sum_{u'\in\pi_e^{-1}(\bar u')}(M_{pt})_{u',u}=(M_t)_{\bar u',\pi_e u}$. Therefore
\[
(S_eL_{e+1}f)(\bar u')=\sum_{u}f(u)\sum_{u'\in\pi_e^{-1}(\bar u')}(M_{pt})_{u',u}
=\sum_{\bar u}(M_t)_{\bar u',\bar u}\sum_{u\in\pi_e^{-1}(\bar u)}f(u)=(L_eS_ef)(\bar u').\qedhere
\]
\end{proof}

\begin{theorem}[Splitting]\label{thm:split}
For every $e\ge1$, $Z_{e+1}:=\ker S_e$ satisfies:
\begin{enumerate}
\item $Z_{e+1}=J_e(H_e)^\perp$ exactly;
\item $Z_{e+1}$ is invariant under $L_{e+1,pt,x}$;
\item $H_{e+1}=J_e(H_e)\oplus Z_{e+1}$ is an orthogonal direct sum of two $L_{e+1,pt,x}$-invariant subspaces, and in an orthonormal basis respecting this sum,
\[
L_{e+1,pt,x}\cong L_{e,t,x}\oplus N_{e+1,pt,x},\qquad N_{e+1,pt,x}:=L_{e+1,pt,x}\big|_{Z_{e+1}}.
\]
\end{enumerate}
\end{theorem}

\begin{proof}
(1) $\langle h,J_e\bar g\rangle_{e+1}=\frac{1}{|K_{e+1}|}\sum_{\bar u}(S_eh)(\bar u)\overline{\bar g(\bar u)}$, so $\langle h,J_e\bar g\rangle_{e+1}=0$ for all $\bar g$ iff $\sum_{u\in\pi_e^{-1}(\bar u)}h(u)=0$ for every $\bar u$, i.e.\ $S_eh=0$.
(2) If $S_ef=0$ then $S_e(L_{e+1}f)=L_e(S_ef)=0$ by Lemma~\ref{lem:push}.
(3) Immediate from (1), (2) and Theorem~\ref{thm:intertwine}: two orthogonal subspaces, each invariant and summing to the whole space, give an exactly block-diagonal matrix in any adapted orthonormal basis. Since $J_e$ is an isometry, the block on $J_e(H_e)$ is unitarily equivalent to $L_{e,t,x}$.
\end{proof}

This is strictly stronger than the invariant-subspace statement alone: Theorem~\ref{thm:intertwine} by itself gives only block-upper-triangularity, for which the complementary block need not be invariant. The operators here are not normal, so the invariance of $Z_{e+1}$ is genuinely additional information, supplied by Lemma~\ref{lem:push} (equivalently, $J_e(H_e)$ is invariant under the adjoint as well). We have $\dim Z_{e+1}=(m_e-1)|K_e|$; at a step with $m_e=1$ (Lemma~\ref{lem:cm}) the new sector is trivial and $L_{e+1,pt,x}$ is unitarily equivalent to $L_{e,t,x}$. The intertwining identities underlying Theorem~\ref{thm:split} were verified in exact arithmetic over thousands of steps, including stalled ones (Section~\ref{sec:numerics}).

\begin{corollary}\label{cor:spec}
$\Spec(L_{e+1,pt,x})=\Spec(L_{e,t,x})\sqcup\Spec(N_{e+1,pt,x})$ (multiset union, with full Jordan structure inherited), and
\[
\rho(L_{e+1,pt,x})=\max\big\{\rho(L_{e,t,x}),\,\rho(N_{e+1,pt,x})\big\}\quad\text{exactly}.
\]
\end{corollary}

\begin{corollary}\label{cor:sv}
The singular-value spectrum of $L_{e+1,pt,x}$ is exactly the union of the singular-value spectra of $L_{e,t,x}$ and $N_{e+1,pt,x}$, since the adjoint is block-diagonal in the same orthonormal basis.
\end{corollary}

\section{Monotonicity}

Write $B_p(e):=\max_{t\ne0}\rho(L_{e,t,x_*})$ and $g_0(p,e):=1-B_p(e)/\lambda_0$, where $\lambda_0=\rho(L_{e,0,x_*})$ is constant in $e$.

\begin{theorem}[Monotonicity]\label{thm:mono}
For every prime $p\nmid ab$ and every $e\ge1$, $B_p(e+1)\ge B_p(e)$, hence $g_0(p,e+1)\le g_0(p,e)$.
\end{theorem}

\begin{proof}
Let $t_\star$ achieve $B_p(e)=\rho(L_{e,t_\star,x_*})$ (it exists, the set being finite; $t_\star$ need not be primitive). Since $t_\star\not\equiv0\pmod{p^e}$, the lift $pt_\star$ is nonzero mod $p^{e+1}$. By Corollary~\ref{cor:spec},
$\rho(L_{e+1,pt_\star,x_*})=\max\{B_p(e),\rho(N_{e+1,pt_\star,x_*})\}\ge B_p(e)$. Since $B_p(e+1)$ is a maximum over all nonzero residues, including $pt_\star$, we get $B_p(e+1)\ge B_p(e)$.
\end{proof}

\begin{corollary}\label{cor:conv}
For every prime $p\nmid ab$, $(g_0(p,e))_{e\ge1}$ is a non-increasing sequence in $[0,1]$, hence convergent, with no exceptional initial levels. Nonnegativity: $\rho(L_{e,t,x_*})\le\rho(L_{e,0,x_*})=\lambda_0$ for every $t$, since $|M_t|=M_0$ entrywise.
\end{corollary}

\begin{remark}
This does not determine the limit. Two logically separate questions remain open: whether only finitely many levels $e$ set a new record for $B_p$ (equivalently, whether the limit is attained at a finite level), and whether the limit is strictly less than $\lambda_0$ (mixing survives) or equal to it ($g_0\to0$). Theorem~\ref{thm:mono} answers neither; it establishes only that the sequence is monotone and therefore has a limit.
\end{remark}

The proof of Theorem~\ref{thm:mono} isolates exactly what remains open:

\begin{openproblem}[Stratum-local dominance]\label{op:local}
For a non-primitive $t$ (writing $t=pt'$), is
\[\rho(N_{e+1,pt',x_*})\le\rho(L_{e,t',x_*}),\]
i.e.\ does the new sector ever exceed its own inherited value? This is not needed for Theorem~\ref{thm:mono}, which requires only one safe witness (the lift of the historical maximizer), but it is the natural finer question about the internal structure of the tower, and it is not resolved here.
\end{openproblem}

\section{Explicit model for the new sector}

$Z_{e+1}=\ker S_e$ is concretely the space of functions on $K_{e+1}$ with zero sum on every fiber of $\pi_e$. Since $K_{e+1}$ is cyclic, fixing a generator $g$ gives a canonical character basis $\chi_k$ ($k=0,\dots,|K_{e+1}|-1$, via the discrete logarithm); $J_e(H_e)$ is spanned exactly by the $\chi_k$ trivial on $\ker\pi_e$, and $Z_{e+1}$ by the remaining ones. When $K_{e+1}$ is the full unit group (primes for which $\langle a,b\rangle$ has index $1$), this coincides with the classical characters of exact conductor $p^{e+1}$; when $K_{e+1}$ is a proper subgroup, the identification with primitive Dirichlet characters is not automatic and is not claimed.

The untwisted operator $L_{e+1,0,x}$ is a convolution on the cyclic group $K_{e+1}$, so the $\chi_k$ diagonalize it; they do not diagonalize $L_{e+1,t,x}$ for $t\ne0$ in general, and no closed form for the eigenvalues of $N_{e+1,t,x}$ at $t\ne0$ was found. In general $\dim Z_{e+1}=(m_e-1)|K_e|$; whenever the group grows ($m_e=p$, which includes every step at which the order of $b$ lifts) this is a constant fraction $(p-1)/p$ of the total dimension, neither vanishing nor dominating.

\section{Arithmetic anomalies and the increment ladder}\label{sec:arith}

The theorems of Sections~3--4 hold at every step. Arithmetic enters only through the \emph{shape} of the increment ladder, that is, through the value of $c_e$ (Lemma~\ref{lem:ladder} and Remark~\ref{rem:stall}). This section records exactly when the order of $b$ stalls.

\begin{proposition}\label{prop:wieferich}
For $p\nmid c$: $c^{p-1}\equiv1\pmod{p^2}$ if and only if $\ord_{p^2}(c)=\ord_p(c)$.
\end{proposition}

\begin{proof}
$\ord_{p^2}(c)\in\{\ord_p(c),\,p\cdot\ord_p(c)\}$ always (standard order lifting for prime powers~\cite{IrelandRosen}). Now $\ord_p(c)\mid p-1$ by Fermat, while $p\cdot\ord_p(c)\ge p>p-1$ cannot divide $p-1$. So $c^{p-1}\equiv1\pmod{p^2}\iff\ord_{p^2}(c)\mid p-1\iff\ord_{p^2}(c)=\ord_p(c)$.
\end{proof}

The following standard consequence of the lifting-the-exponent lemma gives the exact order-lifting schedule.

\begin{proposition}[Order-lifting schedule]\label{prop:lte}
Let $p$ be an odd prime, $p\nmid b$, $r_1=\ord_p(b)$, and $w=w_p(b):=v_p(b^{r_1}-1)\ge1$ (equivalently $w=v_p(b^{p-1}-1)$). Then
\[
\ord_{p^e}(b)=\begin{cases} r_1, & e\le w,\\ r_1\,p^{\,e-w}, & e\ge w.\end{cases}
\]
In particular $c_e=1$ for $1\le e<w$ and $c_e=p$ for $e\ge w$. For non-Wieferich primes to base $b$ ($w=1$) the ladder grows at every step. The depth $w_p(b)$ is thus a parameter of the microscopic ladder geometry; it is not a threshold for inheritance, which holds at every step by Theorem~\ref{thm:split}.
\end{proposition}

\begin{proof}
For odd $p$ and $p\mid c-1$, the lifting-the-exponent lemma gives $v_p(c^n-1)=v_p(c-1)+v_p(n)$. Since $b^n\equiv1\pmod p$ forces $r_1\mid n$, apply this with $c=b^{r_1}$: $v_p(b^{r_1n'}-1)=w+v_p(n')$. Thus $b^{r_1n'}\equiv1\pmod{p^e}$ iff $v_p(n')\ge e-w$, and the smallest such $n'$ is $p^{\max(0,e-w)}$. The equality $v_p(b^{r_1}-1)=v_p(b^{p-1}-1)$ follows from the same formula, since $(p-1)/r_1$ is prime to $p$.
\end{proof}

\paragraph{$p=11$ (Wieferich base 3).}
$3^{10}\equiv1\pmod{121}$, verified directly; $11$ is one of the two known base-3 Wieferich primes ($11$ and $1006003$), and the search of Dorais and Klyve over bases $3,5,7$ found no further base-3 solutions up to about $9.7\times10^{14}$~\cite{Montgomery1993,DoraisKlyve2011}. Here $w_{11}(2)=1$, so the ladder of the resolution base $2$ grows at every level ($c_e=11$). The anomaly sits in the multiplier channel: $\ord_{11^e}(3)=5,5,55,605$ for $e=1,\dots,4$. It may influence the internal spectral data of the new-sector operators, and the previously observed three-level transient in $P_{11}(e)$, followed by a plateau at the next computed level, now sits inside a provably monotone global sequence; but no causal link between $w_{11}(3)=2$ and specific spectral values is claimed here.

\paragraph{$p=1093$ (Wieferich base 2).}
$2^{1092}\equiv1\pmod{1093^2}$, verified directly; $1093$ is the smaller of the two known base-2 Wieferich primes, $1093$ and $3511$~\cite{Wieferich1909,DoraisKlyve2011}. Direct computation gives $\ord_{1093}(2)=\ord_{1093^2}(2)=364$ and $\ord_{1093^3}(2)=364\times1093$, so the ladder stalls exactly at $e=1\to2$ ($c_1=1$), in accordance with Proposition~\ref{prop:lte} ($w=2$). The full group $K_2=\langle 2,3\rangle$ nevertheless grows by a factor $1093$ although the order of $2$ is unchanged: $|K_1|=364$ and $|K_2|=397{,}852$, so $m_1=1093$. This is the case $(c_e,m_e)=(1,p)$ of Lemma~\ref{lem:cm}, with $\langle 2\rangle\cap\ker\pi_1=\{1\}$. The kernel can be named exactly: $3^7=2187=2\cdot1093+1$, so $3^7\equiv1\pmod{1093}$ but $3^7\not\equiv1\pmod{1093^2}$; thus $3^7$ is a nontrivial element of $\ker\pi_1$, which has prime order $1093$, and therefore $\ker\pi_1=\langle 3^7\rangle$. By Lemma~\ref{lem:ladder} each target then receives exactly one fine increment from each reachable fiber, carrying exactly the coarse weight (Remark~\ref{rem:stall}), and the splitting and monotonicity hold at this step as at every other. This was confirmed by an exact entrywise check at sampled entries (Section~\ref{sec:numerics}).

\paragraph{Doubly anomalous primes.}
A prime anomalous in both channels of the Collatz pair would need $w_p(2)\ge2$ and $w_p(3)\ge2$. The only base-2 Wieferich primes below $6.7\times10^{15}$ are $1093$ and $3511$~\cite{DoraisKlyve2011}. Direct computation gives $w_{1093}(3)=w_{3511}(3)=1$, so no doubly anomalous prime exists below $6.7\times10^{15}$. This is a search bound, not an impossibility statement. By Theorem~\ref{thm:split}, none would affect inheritance in any case.

\section{Numerical verification}\label{sec:numerics}

Computation supports the proofs above; it does not replace them.

\paragraph{Exact intertwining audit.}
Both identities $L_{e+1,pt,x}J_e=J_eL_{e,t,x}$ and $S_eL_{e+1,pt,x}=L_{e,t,x}S_e$ were checked entry by entry in exact rational arithmetic ($x=3/8$, weights as fractions, phases kept as exact residues). The sweep covered every prime $p\in\{3,5,7,11,13\}$, every pair $2\le a\ne b\le25$ with $p\nmid ab$, and steps $e\to e+1$ up to $e=2$ (up to $e=3$ for $p\le5$), with group order at most $400$: 3{,}134 steps in all, and zero mismatches. Of these steps, $2{,}872$ were of type $(c_e,m_e)=(p,p)$, $234$ of type $(1,p)$ and $28$ of type $(1,1)$; type $(p,1)$ never occurred, in accordance with Lemma~\ref{lem:cm}. In every case the source-set sizes $|A(u',\bar u)|$ equalled $c_e$, as Lemma~\ref{lem:ladder} predicts.

\paragraph{Targeted controls.}
The following were checked exactly in all entries, and additionally block-diagonalized numerically in an adapted orthonormal basis:
\begin{itemize}
\item generic Collatz, $(3,2)$ at $p=13$, $e=1\to2$;
\item role swap, $(2,3)$ at $p=11$, $e=1\to2$ (stall) and $e=2\to3$;
\item a synthetic depth-3 control, $b=7^5$ with $a=2$ at $p=5$, where $w_5(b)=3$: steps $1\to2$ and $2\to3$ stall, and $3\to4$ is generic; checked with primitive and non-primitive $t$;
\item the multiplier-5 control, $(5,2)$ at $p=7$;
\item a double stall, $(7,18)$ at $p=5$, of type $(1,1)$.
\end{itemize}
Off-diagonal blocks were at most $2\times10^{-15}$, and $\rho(L_{\mathrm{hi}})=\max\{\rho(L_{\mathrm{lo}}),\rho(N)\}$ held to within $3\times10^{-15}$.

\paragraph{The case $p=1093$.}
For the Collatz pair at $p=1093$, $e=1\to2$, the full matrix is too large to build. Instead, 900 entries of the two identities were sampled and checked exactly; there were zero mismatches, with $|A|=1$ throughout.

\paragraph{Earlier checks.}
The earlier verification (exact \texttt{sympy} and high-precision linear algebra at $p=5,13$, including the singular-value union of Corollary~\ref{cor:sv}, matched to $10^{-8}$) is consistent with all of the above. The Wieferich computations of Section~\ref{sec:arith} were made by direct modular exponentiation.

\paragraph{Beyond odd prime powers.}
The algebra underlying the intertwining identities suggests a broader scope: it uses principally the compatibility of reduction with the group structure and the divisibility $r_e\mid r_{e+1}$, whereas other parts of the note use more (Lemma~\ref{lem:cm} uses that $\ker\pi_e$ has order $1$ or $p$, and Section~5 uses cyclicity of $K_e$). Exact checks on 26 additional examples that are not odd prime powers (towers of powers of $2$ with odd $(a,b)$, and composite towers $N\mid N'$ with the twist lifted by $N'/N$) support such an extension. We do not formulate or prove that broader generalization here.

The verification scripts accompany this deposit in the \texttt{src/} directory.

\section{Conclusion}

The prime-power resolution tower of accelerated-Collatz-carry-sum transfer operators splits exactly and orthogonally at every step, for every prime $p\nmid ab$, with no order-lifting hypothesis: order stalls of the resolution base, including the base-2 Wieferich stall at $p=1093$, change only the microscopic shape of the increment ladder, not the inheritance. This upgrades a previously conditional monotonicity claim for the equidistribution spectral gap $g_0(p,e)$ to an unconditional theorem valid from $e=1$, and sharpens the remaining open question to a precise, stratum-local statement (Open Problem~\ref{op:local}) that is no longer entangled with the coarser question the research program originally posed. The theorem and its proof are generic to affine-cocycle transfer operators of this shape; the Collatz carry sum is the motivating instance, not a hypothesis of any result proved here.

\section*{AI Use Disclosure}
AI language-model assistants, including Claude (Anthropic), were used in preparing this note: to help draft and edit the text, typeset it in \LaTeX, review the proofs, suggest standard literature references, independently re-implement the numerical checks of Section~\ref{sec:numerics}, and audit the proofs, which led to the removal of an order-lifting hypothesis used in an earlier draft. All mathematical statements, proofs, computations and references were checked by the author, who takes full responsibility for the content. AI systems are not listed as authors.

\section*{License}
This work is licensed under the Creative Commons Attribution 4.0 International License (CC BY 4.0), \url{https://creativecommons.org/licenses/by/4.0/}.

\begin{thebibliography}{99}

\bibitem{DJ-fixedprime}
E.~De Jes\'us,
\emph{Fixed-Prime Equidistribution of Accelerated Collatz Carry Sums: A Transfer-Operator Spectral Gap, with Two Open Uniformity Problems},
Zenodo, 2026.
\doi{10.5281/zenodo.20694734}.

\bibitem{DJ-amgap}
E.~De Jes\'us,
\emph{An Additive--Multiplicative Spectral-Gap Analyzer for the Fixed-Prime Collatz Carry-Sum Transfer Operator: Instrument, Validation Gates, and Two Exploratory Findings},
technical note, Zenodo, June 2026.

\bibitem{DJ-css}
E.~De Jes\'us,
\emph{Closure, Separation, Stability: A Three-Level Audit Protocol for Reduced Representations},
Zenodo, 2026.
\doi{10.5281/zenodo.21855642}.

\bibitem{Baladi2000}
V.~Baladi,
\emph{Positive Transfer Operators and Decay of Correlations},
Advanced Series in Nonlinear Dynamics, vol.~16, World Scientific, Singapore, 2000.

\bibitem{DoraisKlyve2011}
F.~G.~Dorais and D.~Klyve,
A Wieferich prime search up to $6.7\times10^{15}$,
\emph{Journal of Integer Sequences} \textbf{14} (2011), Article 11.9.2.

\bibitem{IrelandRosen}
K.~Ireland and M.~Rosen,
\emph{A Classical Introduction to Modern Number Theory}, 2nd ed.,
Graduate Texts in Mathematics, vol.~84, Springer, New York, 1990.

\bibitem{Lagarias1985}
J.~C.~Lagarias,
The $3x+1$ problem and its generalizations,
\emph{American Mathematical Monthly} \textbf{92} (1985), no.~1, 3--23.

\bibitem{Lagarias2010}
J.~C.~Lagarias (ed.),
\emph{The Ultimate Challenge: The $3x+1$ Problem},
American Mathematical Society, Providence, RI, 2010.

\bibitem{Montgomery1993}
P.~L.~Montgomery,
New solutions of $a^{p-1}\equiv1\pmod{p^2}$,
\emph{Mathematics of Computation} \textbf{61} (1993), no.~203, 361--363.

\bibitem{Tao2022}
T.~Tao,
Almost all orbits of the Collatz map attain almost bounded values,
\emph{Forum of Mathematics, Pi} \textbf{10} (2022), e12.

\bibitem{Wieferich1909}
A.~Wieferich,
Zum letzten Fermat'schen Theorem,
\emph{Journal f\"ur die reine und angewandte Mathematik} \textbf{136} (1909), 293--302.

\end{thebibliography}
\end{document}
